\subsection{Homomorphisms and Isomorphisms}

\begin{exercise}[1.6.1]
    Let $\phi : G \to H$ be a homomorphism.
    \begin{subexercises}
        \item Prove that $\phi(x^n) = \phi(x)^n$ for all $n \in \zp$.
        \item Do part (a) for $n = - 1$ and deduce that $\phi(x^n) = \phi(x)^n$ for all $n \in \z$.
    \end{subexercises}
\end{exercise}

\begin{solenum}
    \item Inducting on $n$, the base case of $n = 1$ is trivial. If the result holds for $n \in \zp$, then
    \[\phi(x^{n + 1}) = \phi(x^nx) = \phi(x^n)\phi(x) = \phi(x)^n\phi(x) = \phi(x)^{n + 1}.\]
    \item Let $1_G$ and $1_H$ be the identity elements of $G$ and $H$, respectively. Then for $n = 0$, we have
    \[\phi(1_G) = \phi(1_G \cdot 1_G) = \phi(1_G)\phi(1_G),\]
    so that $\phi(1_G) = 1_H$. For $n = -1$ and $x \in G$, we have
    \[1_H = \phi(1_G) = \phi(xx\inv) = \phi(x)\phi(x\inv)\]
    so that $\phi(x\inv) = \phi(x)\inv$. Then for $n < -1$, we have
    \[\phi(x^n) = \phi((x\inv)^{-n}) = \phi(x\inv)^{-n} = \phi(x)^{-n} = \phi(x)^n. \qh\]
\end{solenum}

\begin{specialexercise}[1.6.2]
    If $\phi : G \to H$ is an isomorphism, prove that $|\phi(x)| = |x|$ for all $x \in G$. Deduce that any two isomorphic groups have the same number of elements of order $n$ for each $n \in \z^+$. Is the result true if $\phi$ is only assumed to be a homomorphism?
\end{specialexercise}

\begin{sol}
    Let $x \in G$. Let $\abs x = m$ and $\abs{\phi(x)} = n$. Then
    \[1_H = \phi(1_G) = \phi(x^m) = \phi(x)^m\]
    so that $n \leq m$. Similarly,
    \[1_G = \phi\inv(1_H) = \phi\inv(\phi(x)^n) = x^n\]
    so that $m \leq n$. It follows that $m = n$, i.e., $|\phi(x)| = |x|$. In particular, this shows that both $x$ and $\phi(x)$ both have either finite or infinite order. In particular, if either $|x|$ or $|\phi(x)|$ had infinite order and the other finite, then this would contradict both having finite order. Moreover, since $\phi$ is bijective, it must be bijective between the elements of order $n$ in $G$ and the elements of order $n$ in $H$. Hence, any two isomorphic groups have the same number of elements of order $n$ for each $n \in \z^+$.

    The result is not true if $\phi$ is only assumed to be a homomorphism: the map $\phi : G \to 1$ is a homomorphism for any group $G$, but $G$ may have elements of various orders while $1$ has only the identity element of order 1.
\end{sol}

\begin{exercise}[1.6.3]
    If $\phi : G \to H$ is an isomorphism, prove that $G$ is Abelian if and only if $H$ is Abelian. If $\phi : G \to H$ is a homomorphism, what additional conditions on $\phi$ (if any) are sufficient to ensure that if $G$ is Abelian, then so is $H$?
\end{exercise}

\begin{sol}
    Note that $\phi$ is bijective, hence $\phi\inv$ exists and is also an isomorphism.
    \begin{itemize}
        \item [\rightimp] Suppose that $G$ is Abelian. Then for $h_1 = \phi(g_1), h_2 = \phi(g_2) \in H$ for some $g_1, g_2 \in G$, we have
        \[h_1h_2 = \phi(g_1)\phi(g_2) = \phi(g_1g_2) = \phi(g_2g_1) = \phi(g_2)\phi(g_1) = h_2h_1.\]
        Hence, $H$ is Abelian.
        \item [\leftimp] Suppose that $H$ is Abelian. Then for $g_1 = \phi\inv(h_1), g_2 = \phi\inv(h_2) \in G$ for some $h_1, h_2 \in H$, we have
        \[g_1g_2 = \phi\inv(h_1)\phi\inv(h_2) = \phi\inv(h_1h_2) = \phi\inv(h_2h_1) = \phi\inv(h_2)\phi\inv(h_1) = g_2g_1.\]
        Hence, $G$ is Abelian.
    \end{itemize}
    Let $\phi : G \to H$ be a homomorphism. In the \rightimp direction, observe that the only property of $\phi$ used was its surjectivity to guarantee the existence of $g_1$ and $g_2$ such that $\phi(g_1) = h_1$ and $\phi(g_2) = h_2$. Hence, if $\phi$ is surjective, then $H$ is Abelian whenever $G$ is Abelian.
\end{sol}

\begin{exercise}[1.6.4]
    Prove that the multiplicative groups $\r - \{0\}$ and $\mathbb{C} - \{0\}$ are not isomorphic.
\end{exercise}

\begin{sol}
    Note that $i \in \c \nz$ has order 4, but there is no element in $\r \nz$ that has order 4.
\end{sol}

\begin{exercise}[1.6.5]
    Prove that the additive groups $\r$ and $\q$ are not isomorphic.
\end{exercise}

\begin{sol}
    $\r$ is uncountable, while $\q$ is countable. But bijections preserve cardinality.
\end{sol}

\begin{exercise}[1.6.6]
    Prove that the additive groups $\z$ and $\q$ are not isomorphic.
\end{exercise}

\begin{sol}
    Every nonzero element of $\z$ has infinite order. However, $1/k \in \q$ has order $k$.
\end{sol}

\begin{exercise}[1.6.7]
    Prove that $D_8$ and $Q_8$ are not isomorphic.
\end{exercise}

\begin{sol}
    $D_8$ has only 1 element of order 4, while $Q_8$ has 3 elements of order 4.
\end{sol}

\begin{exercise}[1.6.8]
    Prove that if $n \neq m$, then $S_n$ and $S_m$ are not isomorphic.
\end{exercise}

\begin{sol}
    If $n \neq m$, then $n! \neq m!$ so that $|S_n| \neq |S_m|$.
\end{sol}

\begin{exercise}[1.6.9]
    Prove that $D_{24}$ and $S_4$ are not isomorphic.
\end{exercise}

\begin{sol}
    $D_{24}$ has elements of order 12, but $S_4$ does not.
\end{sol}

\begin{exercise} [1.6.10]
    Fill in the details of the proof that the symmetric groups $S_\Delta$ and $S_\Omega$ are isomorphic if $|\Delta| = |\Omega|$ as follows: let $\theta : \Delta \to \Omega$ be a bijection. Define
    \[
    \phi : S_\Delta \to S_\Omega \quad \text{by} \quad \phi(\sigma) = \theta \circ \sigma \circ \theta^{-1} \quad \text{for all } \sigma \in S_\Delta
    \]
    and prove the following:
    \begin{subexercises}
        \item Prove that $\phi$ is well-defined, that is, if $\sigma$ is a permutation of $\Delta$ then $\theta \circ \sigma \circ \theta^{-1}$ is a permutation of $\Omega$.
        \item Prove that $\phi$ is a bijection from $S_\Delta$ onto $S_\Omega$ by finding a two-sided inverse for $\phi$.
        \item Prove that $\phi$ is a homomorphism, that is, $\phi(\sigma \circ \tau) = \phi(\sigma) \circ \phi(\tau)$.
    \end{subexercises}
    \remark{Note the similarity to the change of basis or similarity transformations for matrices (we shall see the connections between these later in the text)}.
\end{exercise}

\begin{solenum}
    \item Since $\theta : \Delta \to \Omega$ and $\theta\inv : \Omega \to \Delta$ are both bijections with $\sigma$ a bijection, it follows their composition is a bijection. Moreover, $\theta \circ \sigma \circ \theta^{-1}$ is a permutation of $\Omega$, so that $\phi$ is well-defined.
    \item Consider the mapping 
    \[\psi : S_\Omega \to S_\Delta\quad \text{by} \quad \psi(\tau) = \theta\inv \circ \tau \circ \theta \quad \text{for all } \tau \in S_\Omega\]
    By the above discussion, $\psi$ is well-defined. For any $\sigma \in S_\Delta$, we have
    \[\psi(\phi(\sigma)) = \theta\inv \circ (\theta \circ \sigma \circ \theta\inv) \circ \theta = \id_{S_\Delta} \circ \sigma \circ \id_{S_\Delta} = \sigma\]
    and for any $\tau \in S_\Omega$, we have
    \[\phi(\psi(\tau)) = \theta \circ (\theta\inv \circ \tau \circ \theta) \circ \theta\inv = \id_{S_\Omega} \circ \tau \circ \id_{S_\Omega} = \tau\]
    so that $\psi$ is a two-sided inverse for $\phi$. Hence, $\phi$ is a bijection from $S_\Delta$ onto $S_\Omega$.
    \item Let $\sigma, \tau \in S_\Delta$. Then
    \[\phi(\sigma \circ \tau) = \theta \circ (\sigma \circ \tau) \circ \theta\inv = (\theta \circ \sigma) \circ (\tau \circ \theta\inv) = (\theta \circ \sigma \circ \theta\inv) \circ (\theta \circ \tau \circ \theta\inv) = \phi(\sigma) \circ \phi(\tau)\]
    so that $\phi$ is a homomorphism. By the above, $\phi$ is an isomorphism, and hence $S_\Delta \cong S_\Omega$.
\end{solenum}

\begin{exercise}[1.6.11]
    Let $A$ and $B$ be groups. Prove that $A \times B \cong B \times A$.
\end{exercise}

\begin{sol}
    The mapping
    \[\phi : A \times B \to B \times A \quad \text{given by} \quad (a, b) \mapsto (b, a)\]
    for all $a \in A$ and $b \in B$ is clearly bijective. Moreover, for any $(a_1, b_1), (a_2, b_2) \in A \times B$, we have
    \[\phi((a_1, b_1)(a_2, b_2)) = \phi((a_1a_2, b_1b_2)) = (b_1b_2, a_1a_2) = (b_1, a_1)(b_2, a_2) = \phi(a_1, b_1)\phi(a_2, b_2)\]
    so that $\phi$ is a homomorphism. Therefore, $\phi$ is an isomorphism and $A \times B \cong B \times A$.
\end{sol}

\begin{exercise}[1.6.12]
    Let $A$, $B$, and $C$ be groups and let $G = A \times B$ and $H = B \times C$. Prove that $G \times C \cong A \times H$.
\end{exercise}

\begin{sol}
    Consider the mapping
    \[\phi : G \times C \to A \times H \quad \text{by} \quad ((a, b), c) \mapsto (a, (b, c))\]
    for all $a \in A$, $b \in B$, and $c \in C$. The proof is similar to the previous exercise, hence we omit the details to conclude that $\phi$ is an isomorphism and $G \times C \cong A \times H$.
\end{sol}

\begin{specialexercise}[1.6.13]
    Let $G$ and $H$ be groups and let $\phi : G \to H$ be a homomorphism. Prove that the image of $\phi$, $\phi(G)$, is a \defref{subgroup} of $H$. Prove that if $\phi$ is injective then $G \cong \phi(G)$.
\end{specialexercise}

\begin{sol}
    Observe that $1_H \in \phi(G)$ since $\phi(1_G) = 1_H$, so that $\phi(G)$ is nonempty. To show closure, let $h_1, h_2 \in \phi(G)$. Then there exist $g_1, g_2 \in G$ such that $\phi(g_1) = h_1$ and $\phi(g_2) = h_2$. Then
    \[h_1h_2 = \phi(g_1)\phi(g_2) = \phi(g_1g_2) \in \phi(G)\]
    so that $h_1h_2 \in \phi(G)$. Finally,
    \[h_1\inv = \phi(g_1)\inv = \phi(g_1\inv) \in \phi(G)\]
    so that $h_1\inv \in \phi(G)$. Therefore, $\phi(G)$ is a subgroup of $H$.

    Suppose $\phi$ is injective. Then the mapping
    \[\psi : G \to \phi(G) \quad \text{by} \quad g \mapsto \phi(g)\]
    is well-defined. For any $g_1, g_2 \in G$,
    \[\psi(g_1g_2) = \phi(g_1g_2) = \phi(g_1)\phi(g_2) = \psi(g_1)\psi(g_2)\]
    so that $\psi$ is a homomorphism. Moreover, if $\psi(g_1) = \psi(g_2)$, then $\phi(g_1) = \phi(g_2)$, and since $\phi$ is injective, $g_1 = g_2$. Hence, $\psi$ is injective. For any $h \in \phi(G)$, there exists $g \in G$ such that $\phi(g) = h$, and hence $\psi(g) = h$. Therefore, $\psi$ is surjective, hence $\psi$ is an isomorphism and $G \cong \phi(G)$.
\end{sol}

\begin{specialexercise}[1.6.14]
    \label{def:kernel}
    Let $G$ and $H$ be groups and let $\phi : G \to H$ be a homomorphism. Define the \textbf{kernel} of $\phi$ to be $\{ g \in G \mid \phi(g) = 1_H \}$ (so the kernel is the set of elements in $G$ which map to the identity of $H$, i.e., is the fiber over the identity of $H$). Prove that the kernel of $\phi$ is a subgroup of $G$. Prove that $\phi$ is injective if and only if the kernel of $\phi$ is the identity subgroup of $G$.
\end{specialexercise}

\begin{sol}
    Denote the kernel of $\phi$ by $\ker\phi$. Since $\phi(1_G) = 1_H$, then $1_G \in \ker\phi$, so that $\ker\phi$ is nonempty. To show closure, let $g_1, g_2 \in \ker\phi$. Then
    \[\phi(g_1g_2) = \phi(g_1)\phi(g_2) = 1_H1_H = 1_H\]
    so that $g_1g_2 \in \ker\phi$. Finally,
    \[\phi(g_1\inv) = \phi(g_1)\inv = 1_H\inv = 1_H\]
    so that $g_1\inv \in \ker\phi$. Therefore, $\ker\phi$ is a subgroup of $G$.
    \begin{itemize}
        \item [\rightimp] If $\phi$ is injective, then for $g \in \ker\phi$, we have $\phi(g) = 1_H = \phi(1_G)$. Then $g = 1_G$, so that $\ker\phi = \{1_G\}$.
        \item [\leftimp] If $\ker\phi = \{1_G\}$, then for $g_1, g_2 \in G$, if $\phi(g_1) = \phi(g_2)$, we have $\phi(g_1)\phi(g_2)\inv = 1_H$, so $g_1g_2\inv \in \ker\phi$. Hence, $g_1g_2\inv = 1_G$, and thus $g_1 = g_2$. Therefore, $\phi$ is injective. \qh
    \end{itemize}
\end{sol}

\begin{exercise}[1.6.15]
    Define a map $\pi : \r^2 \to \r$ by $\pi((x,y)) = x$. Prove that $\pi$ is a homomorphism and find the kernel of $\pi$.
\end{exercise}

\begin{sol}
    Let $(x_1, y_1), (x_2, y_2) \in \r^2$. Then
    \[\pi((x_1, y_1) + (x_2, y_2)) = \pi((x_1 + x_2, y_1 + y_2)) = x_1 + x_2 = \pi((x_1, y_1)) + \pi((x_2, y_2))\]
    so that $\pi$ is a homomorphism. Moreover, if $(x, y) \in \ker(\pi)$, then $\pi((x, y)) = x = 0$. Hence,
    \[\ker(\pi) = \{(0, y) \mid y \in \r\} \qh\]
\end{sol}

\begin{exercise}[1.6.16]
    Let $A$ and $B$ be groups and let $G$ be their direct product, $A \times B$. Prove that the maps $\pi_1 : G \to A$ and $\pi_2 : G \to B$ defined by $\pi_1((a,b)) = a$ and $\pi_2((a,b)) = b$ are homomorphisms and find their kernels.
\end{exercise}

\begin{sol}
    The proof is similar to the previous exercise, hence we omit the details to conclude that $\pi_1$ and $\pi_2$ are homomorphisms. Moreover, if $(a, b) \in \ker(\pi_1)$, then $\pi_1((a, b)) = a = 1_A$. If $(a, b) \in \ker(\pi_2)$, then $\pi_2((a, b)) = b = 1_B$. Hence,
    \[\ker(\pi_1) = \{(1_A, b) \mid b \in B\}, \quad \text{and} \quad \ker(\pi_2) = \{(a, 1_B) \mid a \in A\} \qh\]
\end{sol}

\begin{exercise} [1.6.17]
    Let $G$ be any group. Prove that the map from $G$ to itself defined by $g \mapsto g^{-1}$ is a homomorphism if and only if $G$ is Abelian.
\end{exercise}

\begin{solitem}
    \item [\rightimp] Suppose $\phi$ is a homomorphism. Then for $g, h \in G$, we have
    \[(gh)\inv = \phi(gh) = \phi(g)\phi(h) = g\inv h\inv = h\inv g\inv.\]
    Since inverses are unique, then $hg = gh$ so that $G$ is Abelian.
    \item [\leftimp] Suppose $G$ is Abelian. Then for $g, h \in G$, we have
    \[\phi(gh) = (gh)\inv = h\inv g\inv = g\inv h\inv = \phi(g)\phi(h),\]
    so that $\phi$ is a homomorphism.
\end{solitem}

\begin{exercise}[1.6.18]
    Let $G$ be any group. Prove that the map from $G$ to itself defined by $g \mapsto g^2$ is a homomorphism if and only if $G$ is Abelian.
\end{exercise}

\begin{sol}
    The proof is similar to \exref{1.6.17}, hence we omit the details.
\end{sol}

\begin{exercise}[1.6.19]
    Let $G = \{ z \in \mathbb{C} \mid z^n = 1 \text{ for some } n \in \z^+ \}$. Prove that for any fixed integer $k > 1$ the map from $G$ to itself defined by $z \mapsto z^k$ is a surjective homomorphism but is not an isomorphism.
\end{exercise}

\begin{sol}
    Since $G \subseteq \c$, then $G$ is Abelian. Let $\phi : z \to z^k$. Hence, for any $z, w \in G$, we have
    \[\phi(zw) = (zw)^k = z^kw^k = \phi(z)\phi(w)\]
    so that $\phi$ is a homomorphism. To show surjectivity, let $w \in G$ such that $w^n = 1$ for some $n \in \zp$. Note that $w^{1/k} \in G$ because $(z^{1/k})^{nk} = z^n = 1$. Then
    \[\phi(z^{1/k}) = (z^{1/k})^k = z\]
    so that $\phi$ is surjective. Finally, any $z = e^{2\pi i/k} \in G$ is such that $\phi(z) = z^k = 1$ with $z \neq 1$ since $k > 1$. Hence, $\phi$ is not injective and thus not an isomorphism.
\end{sol}

\begin{specialexercise}[1.6.20]
    \label{def:automorphism group}
    Let $G$ be a group and let $\aut(G)$ be the set of all isomorphisms from $G$ onto $G$. Prove that $\aut(G)$ is a group under function composition (called the \textbf{automorphism group} of $G$ and the elements of $\aut(G)$ are called \textbf{automorphisms} of $G$).
\end{specialexercise}

\begin{sol}
    Observe that $\aut(G)$ is associative under function composition and has identity through the identity map $\id : G \to G$ given by $\id(g) = g$. Moreover, $\phi \circ \psi$ for $\phi, \psi \in \aut(G)$ is a bijection. Then for any $g, h \in G$, we have
    \[(\phi \circ \psi)(gh) = \phi(\psi(gh)) = \phi(\psi(g)\psi(h)) = \phi(\psi(g))\phi(\psi(h)) = (\phi \circ \psi)(g)(\phi \circ \psi)(h),\]
    so that $\phi \circ \psi$ is a homomorphism. Hence, $\phi \circ \psi \in \aut(G)$, and $\aut(G)$ is closed under function composition. Finally, for any $\phi \in \aut(G)$, then $\phi\inv$ exists and is a bijection. Then
    \[\phi\inv(gh) = \phi\inv(\phi(\phi\inv(g))\phi(\phi\inv(h))) = \phi\inv(\phi(\phi\inv(g)\phi\inv(h))) = \phi\inv(g)\phi\inv(h),\]
    so that $\phi\inv$ is a homomorphism. Hence, $\phi\inv \in \aut(G)$, and $(\aut(G), \circ)$ is a group.
\end{sol}

\begin{exercise} [1.6.21]
    Prove that for each fixed nonzero $k \in \q$ the map from $\q$ to itself defined by $q \mapsto kq$ is an automorphism of $\q$.
\end{exercise}

\begin{sol}
    Let $\phi : \q \to \q$ be defined by $\phi(q) = kq$ for all $q \in \q$. For any $p, q \in \q$, we have
    \[\phi(p + q) = k(p + q) = kp + kq = \phi(p) + \phi(q),\]
    so that $\phi$ is a homomorphism. Bijectivity is clear, hence $\phi \in \aut(\q)$.
\end{sol}

\begin{exercise}[1.6.22]
    Let $A$ be an Abelian group and fix some $k \in \z$. Prove that the map $a \mapsto a^k$ is a homomorphism from $A$ to itself. If $k = -1$ prove that this homomorphism is an isomorphism (i.e., is an automorphism of $A$).
\end{exercise}

\begin{sol}
    Let $\phi : A \to A$ be defined by $\phi(a) = a^k$ for all $a \in A$. For any $a_1, a_2 \in A$, we have
    \[\phi(a_1a_2) = (a_1a_2)^k = a_1^k a_2^k = \phi(a_1)\phi(a_2)\]
    so that $\phi$ is a homomorphism. For $k = -1$, see \exref{1.6.17} to conclude that $\phi \in \aut(A)$.
\end{sol}

\begin{specialexercise}[1.6.23]
    Let $G$ be a finite group which possesses an automorphism $\sigma$ such that $\sigma(g) = g$ if and only if $g = 1$. If $\sigma^2$ is the identity map from $G$ to $G$, prove that $G$ is Abelian (such an automorphism $\sigma$ is called \textbf{fixed point free} of order 2).

    \hint{Show that every element of $G$ can be written in the form $x^{-1}\sigma(x)$ and apply $\sigma$ to such an expression.}
\end{specialexercise}

\begin{sol}
    Define the map
    \[f : G \to G \quad \text{by} \quad x \mapsto x\inv \sigma(x)\]
    for all $x \in G$. If $f(x) = f(y)$ for $x, y \in G$, then $x\inv \sigma(x) = y\inv \sigma(y)$ so that $yx\inv = \sigma(yx\inv)$. Since $\sigma$ is fixed point free, then $yx\inv = 1$ and $x = y$. Hence, $f$ is injective. Since $G$ is finite, then $f$ is also surjective, and every element of $G$ may be written in the form $x\inv \sigma(x)$ for some $x \in G$.

    Let $g \in G$ be arbitrary. Then there exists $x \in G$ such that $g = x\inv \sigma(x)$. Applying $\sigma$ to both sides, we have
    \[\sigma(g) = \sigma(x\inv \sigma(x)) = \sigma(x\inv) \sigma^2(x) = \sigma(x)\inv x = (x\inv \sigma(x))\inv = g\inv.\]
    Then for any $g, h \in G$, we have
    \[\sigma(gh) = (gh)\inv = h\inv g\inv = \sigma(h) \sigma(g) = \sigma(hg).\]
    Since $\sigma$ is injective, then $gh = hg$ and $G$ is Abelian.
\end{sol}

\begin{exercise}[1.6.24]
    Let $G$ be a finite group and let $x$ and $y$ be distinct elements of order 2 in $G$ that generate $G$. Prove that $G \cong D_{2n}$, where $n = |xy|$.
\end{exercise}

\begin{sol}
    To show that $G$ is isomorphic to $D_{2n}$, we first need to produce an element $t \in G$ of order $n$ along with an element $x \in G$ of order 2 such that $xt = t\inv x$. The clear choice is to let $t = xy$ and keep $x$ as is. Note that $|t| = |xy| = n$ by assumption, and $|x| = 2$ by assumption. Moreover,
    \[xt = x(xy) = (xx)y = y = (xy)\inv x = t\inv x\]
    so that the elements $x, t \in G$ satisfy the relations of $D_{2n}$. Additionally, observe that $y = xt$ so that $t$ and $x$ generate $G$. Following this, it is easy to see that every element of $G$ can be written in the form $x^i t^j$ for $i \in \{0, 1\}$ and $0 \leq j < n$. For $x = 0$, we have the elements in the cyclic subgroup generated by $t$, and for $x = 1$, we have the elements $\set{x, xt, xt^2, \ldots, xt^{n-1}}$. Since $|t| = n$, then these are all distinct elements of $G$. Hence, $|G| = 2n$.

    Since the relations of $D_{2n}$ are satisfied by $x, t \in G$, $G$ is generated by $x$ and $t$, and $|G| = |D_{2n}| = 2n$, it follows that the map
    \[\phi : D_{2n} \to G \quad \text{by} \quad r \mapsto t, s \mapsto x\]
    extends to an isomorphism from $D_{2n}$ onto $G$. Therefore, $G \cong D_{2n}$.
\end{sol}

\begin{exercise}[1.6.25]
    Let $n \in \z^+$, let $r$ and $s$ be the usual generators of $D_{2n}$ and let $\theta = 2\pi / n$.
    \begin{subexercises}
        \item Prove that the matrix 
        \[
        \begin{pmatrix}
        \cos \theta & -\sin \theta \\
        \sin \theta & \cos \theta
        \end{pmatrix}
        \]
        is the matrix of the linear transformation which rotates the $x,y$-plane about the origin in a counterclockwise direction by $\theta$ radians.
        \item Prove that the map $\phi : D_{2n} \to \gl_2(\r)$ defined on generators by
        \[
        \phi(r) = 
        \begin{pmatrix}
            \cos \theta & -\sin \theta \\
            \sin \theta & \cos \theta
        \end{pmatrix}
        \quad \text{and} \quad
        \phi(s) = 
        \begin{pmatrix}
            0 & 1 \\
            1 & 0
        \end{pmatrix}
        \]
        extends to a homomorphism of $D_{2n}$ into $\gl_2(\r)$.
        \item Prove that the homomorphism $\phi$ defined above is injective.
    \end{subexercises}
\end{exercise}

\begin{solenum}
    \item Let $\r^2$ be generated by the standard basis vectors $\b e_1 = (1 \quad 0)^T$ and $\b e_2 = (0 \quad 1)^T$. Let $T$ denote the linear transformation which rotates the $x,y$-plane about the origin in a counterclockwise direction by $\theta$ radians. Then
    \[T(\b e_1) = 
    \begin{pmatrix}
        \cos \theta \\
        \sin \theta
    \end{pmatrix}, \quad \text{and} \quad
    T(\b e_2) = 
    \begin{pmatrix}
        -\sin \theta \\
        \cos \theta
    \end{pmatrix}.
    \]
    Then the matrix of $T$ with respect to the standard basis $\{\b e_1, \b e_2\}$ is exactly the desired matrix.
    \item It is easy to show by induction that for any positive $k$, then $\phi(r)^k$ is given by the following matrix:
    \[\phi(r)^k = \begin{pmatrix}
        \cos (k\theta) & -\sin (k\theta) \\
        \sin (k\theta) & \cos (k\theta)
    \end{pmatrix}\]
    so that $\phi(r)^n = I$. It is fairly straightforward to see that the set of vectors $\set{T(\b e_1), T(\b e_2)}$ is orthonormal, so $\phi(r)$ is an orthogonal matrix, and hence $\phi(r)\inv = \phi(r)^T$. Moreover, we can calculate that $\phi(s)^2 = I$. Lastly, we calculate that
    \[\phi(s)\phi(r) = 
    \begin{pmatrix}
        0 & 1 \\
        1 & 0
    \end{pmatrix}
    \begin{pmatrix}
        \cos \theta & -\sin \theta \\
        \sin \theta & \cos \theta
    \end{pmatrix} = 
    \begin{pmatrix}
        \sin \theta & \cos \theta \\
        \cos \theta & -\sin \theta
    \end{pmatrix} = 
    \begin{pmatrix}
        \cos \theta & \sin \theta \\
        -\sin \theta & \cos \theta
    \end{pmatrix}
    \begin{pmatrix}
        0 & 1 \\
        1 & 0
    \end{pmatrix} = \phi(r)\inv \phi(s)\]
    so that $\phi(s)\phi(r) = \phi(r)\inv \phi(s)$. Since these relations match those of $D_{2n}$, then $\phi$ extends to a homomorphism from $D_{2n}$ to $\gl_2(\r)$.
    \item To show that $\phi$ is injective, we show that $\ker\phi$ is trivial. Let $g \in \ker\phi$.
    \begin{itemize}
        \item If $g = r^k$ for some $0 \leq k < n$, then $\phi(r^k) \neq I$, since $\phi(r)^k$ is a rotation matrix and only $\phi(r)^n = I$.
        \item If $g = sr^k$ for some $0 \leq k < n$, then $\phi(sr^k) = \phi(s)\phi(r)^k$. Note that $\phi(s)$ has determinant $-1$ while $\phi(r)^k$ has determinant $1$ since it is a rotation matrix. Hence, $\phi(sr^k)$ has determinant $-1$ and cannot be the identity matrix.
    \end{itemize}
    Therefore, the only element in $\ker\phi$ is the identity element of $D_{2n}$, so that $\phi$ is injective.
\end{solenum}

\begin{exercise}[1.6.26]
    Let $i$ and $j$ be the generators of $Q_8$ described in Section 5. Prove that the map $\phi$ from $Q_8$ to $\gl_2(\mathbb{C})$ defined on generators by
    \[\phi(i) = 
    \begin{pmatrix}
        \sqrt{-1} & 0 \\
        0 & -\sqrt{-1}
    \end{pmatrix}
    \quad \text{and} \quad
    \phi(j) = 
    \begin{pmatrix}
        0 & -1 \\
        1 & 0
    \end{pmatrix}
    \]
    extends to a homomorphism. Prove that $\phi$ is injective.
\end{exercise}

\begin{sol}
    Matrix multiplication shows that the relations described in \exref{1.5.3} are satisfied, so that $\phi$ extends to a homomorphism from $Q_8$ to $\gl_2(\mathbb{C})$. Moreover, $\phi$ is injective since the only element that maps to the identity matrix is $1 \in Q_8$. Therefore, $\phi$ is an injective homomorphism.
\end{sol}